\documentclass[11pt]{article}
\usepackage[margin=0.9in]{geometry}
\usepackage{amsmath,amssymb,amsthm,booktabs}
\usepackage[T1]{fontenc}
\usepackage{lmodern,microtype}
\usepackage[colorlinks=true,linkcolor=blue,urlcolor=blue]{hyperref}
\newtheorem{theorem}{Theorem}[section]
\newtheorem{lemma}[theorem]{Lemma}
\newtheorem{proposition}[theorem]{Proposition}
\newtheorem{corollary}[theorem]{Corollary}
\theoremstyle{remark}
\newtheorem{remark}[theorem]{Remark}
\newcommand{\norm}[1]{\left\|#1\right\|}
\title{Toward a four-dimensional Yang--Mills construction\\
\large Notebook 13: a uniform nonlinear constraint chart}
\author{Research working document}
\date{16 October 2026}
\begin{document}
\maketitle
\begin{abstract}
The complete corrected averaging map, iterated through fixed-size
blocks, admits a complex neighborhood whose physical radius does
not shrink with refinement. We prove this from its exact local
coordinate, the contraction of its flat linear part and a geometric
sum of physical scale factors. Its differential is uniformly close
to the flat average in bounded-field and physical Hilbert norms.
The flat minimizing extension then gives an exact analytic chart
of the nonlinear constraint, including an exponentially local
right inverse for its ordinary-coordinate differential. These are
fixed-reference statements; the moving-background operator and
the nonlinear action minimizer require separate identification.
\end{abstract}
\noindent\textbf{Status.}
Research proofs checked by separate internal derivations and audits.
No external certification or novelty claim is made. Source matching
is to the pinned deposited transcription identified below.
% Notebook 13 body fragment, no preamble. Unique label prefix: na-.
% Requires amsmath, amssymb, amsthm and a one-argument \norm macro.

\section{A uniform analytic neighborhood for the corrected nonlinear average}
\label{na-section}

Fix an odd one-step ratio $L\ge3$, a refinement $M=L^j$, fine
spacing $\eta=M^{-1}$ and final physical scale one. The reference
background is the identity at every stage. Constants may depend on
$L$, dimension, color norm and the specified averaging prescription,
but not on $j$, $M$ or compatible periodic volume. Pairings and
Hilbert norms use physical lattice volume. All complex analytic maps
below use the fixed flat-reference logarithmic coordinate. Their
differentials at a nonzero field are not identified here with a
separately left-trivialized nonflat source operator.

\subsection{The exact source coordinate and the fixed local constants}
\label{na-source}

Deposited Raw I:24320--24325 composes the corrected nonlinear
operation, with its averaged reference background. Since that
reference is initially the identity, it remains the identity.
I:24344--24353 gives the logarithmic coordinate, its linear part and
the quadratic remainder; I:24357--24369 identifies the flat
linearization. The coordinate bridge I:47076--47087, using the maps
defined at I:46965--46973, gives
\begin{equation}
 F_L(X;c)=\frac1i\log\left[
  e^{-i\alpha_{c_-}(X)}e^{i\phi(X;c)}e^{i\alpha_{c_+}(X)}\right],
 \qquad DF_L(0)=LQ_L^{\mathrm{st}},
 \label{na-coordinate}
\end{equation}
where
\[
 \phi(X;c)=\Phi^{\flat}(1,X;c),\qquad
 \alpha_y(X)=L^{-d}\sum_{x\in B(y)}\Phi^{\mathfrak g}(1,X;y,x).
\]
This identity includes the endpoint corrections; it is not the
uncorrected rooted logarithm alone. These source matches are checked
at the deposited-transcription level. The original CMP98 full paper
has not been independently reverified in this calculation.

For fixed $L$, the products, means and logarithms in
\eqref{na-coordinate} give a finite-dimensional local holomorphic
map near zero. Let $D_L(c)$ be its finite set of input bonds and set
\[
 m_L=\sup_c|D_L(c)|,\qquad
 N_L=\sup_b\#\{c:b\in D_L(c)\}.
\]
All translated local maps share the same constants. Write
$S_L(X)=F_L(X)-LQ_L^{\mathrm{st}}X$. Since
$S_L(0)=DS_L(0)=0$, bounding the second derivative on a smaller
closed complex polydisc gives constants $r_0,C_0,K_0>0$ such that
\begin{equation}
 |S_L(X;c)|\le C_0\norm{X}_\infty^2,
 \qquad
 \norm{\partial_bS_L(X;c)}\le K_0\norm{X}_\infty
 \quad (b\in D_L(c)),
 \label{na-local-jet-bounds}
\end{equation}
whenever $\norm{X}_\infty<r_0$. The derivative is zero outside
$D_L(c)$. On short tori, use the lifted local pattern with
multiplicities. No constant uniform as $L\to\infty$ is required.

\subsection{Physical scaling and the one-step differential}
\label{na-one-step}

At input spacing $h$ and output spacing $Lh$, define
\begin{equation}
 \mathcal F_h(A)=(Lh)^{-1}F_L(hA)=Q_hA+R_h(A),
 \qquad R_h(A)=(Lh)^{-1}S_L(hA).
 \label{na-physical-map}
\end{equation}
Put
\[
 C=C_0/L,\qquad D_\infty=m_LK_0/L,\qquad
 D_2=K_0\sqrt{m_LN_L}\,L^{d/2-1}.
\]
For $h\norm{A}_\infty<r_0$, one has
\begin{align}
 \norm{R_h(A)}_\infty&\le Ch\norm{A}_\infty^2,
 \label{na-one-step-rem}\\
 \norm{DR_h(A)}_{\infty\to\infty}
    &\le D_\infty h\norm{A}_\infty,
 \qquad
 \norm{DR_h(A)}_{2,h\to2,Lh}
    \le D_2h\norm{A}_\infty.
 \label{na-one-step-der}
\end{align}
Indeed each derivative matrix entry is at most
$K_0h\norm{A}_\infty/L$. Its row sum is at most
$m_LK_0h\norm{A}_\infty/L$, and its column sum in the physical
adjoint normalization is at most
$L^dN_LK_0h\norm{A}_\infty/L$. Their geometric mean proves the
Hilbert bound. This is a measure-normalized Schur estimate, with no
factor from the number of fine sites in a final unit cell.

The flat average is
\[
 (Q_hA)_\mu(y)=L^{-d-1}\sum_{x\in B(y)}\sum_{r=0}^{L-1}
                         A_\mu(x+rh e_\mu).
\]
Its coefficients are nonnegative, its row sum is one and its
unweighted column sum is $L^{-d}$: a given fine bond has exactly
$L$ choices of the string offset. The physical adjoint supplies the
factor $L^d$. Thus $Q_h$ and its physical adjoint have row sums one;
in particular $Q_h$ contracts both the bounded-field and physical
Hilbert norms. Multiplicity counting gives the same assertion on
short compatible tori.

\subsection{The complete nonlinear iterate on a fixed complex ball}
\label{na-uniform-ball}

Set $h_t=L^{t-j}$ for $0\le t\le j$ and $S=(L-1)^{-1}$. Let
\[
 A_0=A,\qquad A_{t+1}=\mathcal F_{h_t}(A_t),\qquad
 \mathcal F^{(j)}(A)=A_j,
 \qquad Q_{r:s}=Q_{s-1}\cdots Q_r.
\]
An empty product is the identity. Then $Q=Q_{0:j}=Q_M^{\mathrm{st}}$.

\begin{theorem}[Uniform nonlinear ball]
\label{na-ball-theorem}
Choose
\begin{equation}
 0<\rho<\min\{r_0/4,\,(4CS)^{-1}\},
 \label{na-radius}
\end{equation}
omitting the second restriction when $C=0$. The complete corrected
iterate is holomorphic on $\norm{A}_\infty<\rho$, uniformly in
refinement and volume. If $a=\norm{A}_\infty$, then
\begin{equation}
 \norm{A_t}_\infty\le
 \frac{a}{1-Ca\sum_{s<t}h_s}\le2a.
 \label{na-riccati}
\end{equation}
Writing $R^{(j)}=\mathcal F^{(j)}-Q$ and $C_R=4CS$, one also has
\begin{equation}
 \norm{\mathcal F^{(j)}(A)}_\infty\le2a,
 \qquad \norm{R^{(j)}(A)}_\infty\le C_Ra^2.
 \label{na-composite-remainder}
\end{equation}
\end{theorem}

\begin{proof}
Let $u_t=a/(1-Ca\sum_{s<t}h_s)$. The denominators are positive.
If $\norm{A_t}_\infty\le u_t$, the local analytic domain holds
because $h_tu_t\le2\rho<r_0$. Moreover
\[
 u_t+Ch_tu_t^2\le\frac{u_t}{1-Ch_tu_t}=u_{t+1}.
\]
Induction proves both the estimates and the validity of every local
evaluation, without presupposing the domain bootstrap. Each
composition is therefore holomorphic with a common strict margin.
The exact telescoping identity
\[
 R^{(j)}(A)=\sum_{t=0}^{j-1}Q_{t+1:j}R_{h_t}(A_t)
\]
and the contractions of $Q_{t+1:j}$ give
$\norm{R^{(j)}(A)}_\infty\le4Ca^2\sum_t h_t\le4CSa^2$.
For $j=0$ the map is the identity and the remainder vanishes.
The same local proof applies to the bounded-field Banach space on
the infinite lattice.
\end{proof}

\subsection{Differential estimates and factorial bounds}
\label{na-differentials}

For $p=2,\infty$, the physical intermediate norms when $p=2$ and
the ordered chain rule give
\begin{align}
 \norm{D\mathcal F^{(j)}(A)}_{p\to p}
    &\le\prod_t(1+2D_pa h_t)\le e^{2D_pSa},
 \label{na-derivative-bound}\\
 \norm{D\mathcal F^{(j)}(A)-Q}_{p\to p}
    &\le e^{2D_pSa}-1=:\delta_p(a).
 \label{na-differential-deviation}
\end{align}
For the second line, expand the difference between the product of
$Q_t+DR_{h_t}(A_t)$ and the product of $Q_t$, keeping their
noncommutative order. The nonempty selections of remainder factors
sum in norm to at most $\prod_t(1+2D_pa h_t)-1$. In particular
no nonflat contraction is assumed.

The differential is holomorphic as an operator from fine physical
$L^2$ to coarse physical $L^2$, with the field variable in bounded
norm. For fixed $j$, the local translated Taylor expansions have
finite incidence and converge in this operator norm; the preceding
uniform bounds make the resulting estimates independent of volume.
For $0<r<\rho-\norm{A}_\infty$ and $n\ge1$, Cauchy gives
\begin{equation}
 \norm{D^n\mathcal F^{(j)}(A)[K_1,\ldots,K_n]}_\infty
 \le2\rho\,n!(e/r)^n\prod_i\norm{K_i}_\infty.
 \label{na-cauchy-infty}
\end{equation}
For $n\ge0$, with $n=0$ read directly from
\eqref{na-derivative-bound},
\begin{equation}
 \norm{D^{n+1}\mathcal F^{(j)}(A)[K_1,\ldots,K_n,B]}_{2,1}
 \le n!(e/r)^n e^{2D_2S\rho}
       \norm{B}_{2,\eta}\prod_i\norm{K_i}_\infty.
 \label{na-cauchy-mixed}
\end{equation}
For $n\ge1$ use product-polydisc radii $r/n$ for normalized
directions, followed by $n^n\le e^n n!$. The mixed estimate has one
Hilbert input; all other directions are bounded fields.

\subsection{Finite range and weighted physical-cell operators}
\label{na-weighted}

A one-step reading radius at most $r_Lh$ gives a total physical
reading radius at most $r_L\sum_t h_t\le r_LS$. Thus an integer
unit-cell range $R_*$, including fixed cell-diameter margins,
contains the support of both $D\mathcal F^{(j)}(A)$ and $Q$,
independently of $j$. The sharper deposited two-block dependence
I:24332--24336 is more than needed here.

For fine or coarse physical-cell Hilbert fibers use
\begin{equation}
 \norm{T}_\kappa=\max\left\{
  \sup_x\sum_y e^{\kappa|x-y|_1}\norm{T_{xy}},\,
  \sup_y\sum_x e^{\kappa|x-y|_1}\norm{T_{xy}}\right\}.
 \label{na-schur-norm}
\end{equation}
Let $v(R_*)$ count cells in that range and set
$\chi_\kappa=v(R_*)e^{\kappa R_*}$. Every cell block is bounded
by the global Hilbert operator norm, so
\begin{equation}
 \norm{DR^{(j)}(A)}_\kappa\le\chi_\kappa\delta_2(a).
 \label{na-weighted-deviation}
\end{equation}
All coefficient derivatives retain the same finite range. Hence the
operator-valued Taylor series converge in the weighted norm too,
and $A\mapsto DR^{(j)}(A)$ is holomorphic into this algebra.
The analogous finite-range conversion for bounded-cell fibers uses
$\delta_\infty(a)$. Periodic restrictions have the same bounds by
lifted multiplicities or periodic cell distance. There is no fine
dimension factor and no assertion of a special ordered insertion
path.

\subsection{An exact analytic chart for the nonlinear constraint}
\label{na-chart}

Let $H$ be the already proved source-compatible flat minimizing
extension, with
\begin{equation}
 QH=I,\qquad \norm{H}_{\infty\to\infty}\le k,
 \qquad \norm{H}_\kappa\le k_\kappa.
 \label{na-H-input}
\end{equation}
These constants are uniform in refinement and volume. Choose
$0<r_A<\rho$ such that
\begin{equation}
 kC_Rr_A\le\tfrac14,\qquad
 k\delta_\infty(r_A)\le\tfrac12,\qquad
 k_\kappa\chi_\kappa\delta_2(r_A)\le\tfrac12.
 \label{na-chart-radius}
\end{equation}

\begin{theorem}[Uniform constraint chart]
\label{na-chart-theorem}
For $\xi\in\ker Q$, coarse field $B$, and
$u=\xi+HB$ with $\norm{u}_\infty<r_A/2$, the equation
\begin{equation}
 A=u-HR^{(j)}(A)
 \label{na-fixed-point}
\end{equation}
has a unique solution in $\norm{A}_\infty\le r_A$. This defines
a holomorphic map $A=\mathcal A_j(\xi,B)$ satisfying
\begin{align}
 \mathcal F^{(j)}(\mathcal A_j(\xi,B))&=B,
 & (I-HQ)\mathcal A_j(\xi,B)&=\xi,
 \label{na-chart-identities}\\
 \norm{\mathcal A_j(\xi,B)-\xi-HB}_\infty
   &\le4kC_R\norm{\xi+HB}_\infty^2.
 \label{na-chart-remainder}
\end{align}
Every field with $\norm{A}_\infty<r_A/3$ belongs to this chart.
Its ordinary-coordinate constraint differential has a holomorphic
right inverse with uniform exponential physical-cell locality.
\end{theorem}

\begin{proof}
The right side of \eqref{na-fixed-point} maps the closed radius
$r_A$ ball into itself: its nonlinear term has norm at most
$kC_Rr_A^2\le r_A/4$. Its Lipschitz constant is at most one half
by \eqref{na-differential-deviation} and \eqref{na-chart-radius}.
The contraction theorem gives uniqueness and existence. Uniform
iteration on smaller parameter balls gives holomorphic dependence.

One precise parameter norm is
$\norm{(\xi,B)}_{\mathrm{par}}=\norm{\xi+HB}_\infty$.
It is a norm because $Qu=B$, and it is uniformly equivalent to
the usual direct-sum norm:
\[
 \norm{B}_\infty\le\norm{u}_\infty,\qquad
 \norm{\xi}_\infty\le(1+k)\norm{u}_\infty,\qquad
 \norm{u}_\infty\le\norm{\xi}_\infty+k\norm{B}_\infty.
\]
Apply $Q$ and $I-HQ$ to \eqref{na-fixed-point} to obtain
\eqref{na-chart-identities}. Also
$\norm{A}_\infty\le2\norm{u}_\infty$, since
$\norm{HR^{(j)}(A)}_\infty\le\frac12\norm{A}_\infty$ on this
ball. This proves \eqref{na-chart-remainder}. Conversely set
$\xi=(I-HQ)A$, $B=\mathcal F^{(j)}(A)$. Then
$u=A+HR^{(j)}(A)$ has norm below $r_A/2$ whenever
$\norm{A}_\infty<r_A/3$, so uniqueness identifies its fixed point
with $A$.

At a chart point set
$Q_A=D\mathcal F^{(j)}(A)=Q+DR^{(j)}(A)$. The ordered Neumann
series give
\begin{equation}
 H_A=(I+H DR^{(j)}(A))^{-1}H
      =H(I+DR^{(j)}(A)H)^{-1},\qquad Q_AH_A=I.
 \label{na-chart-right-inverse}
\end{equation}
In the weighted algebra the same series yield
\begin{equation}
 \norm{H_A}_\kappa\le2k_\kappa,
 \qquad \norm{(I+H DR^{(j)}(A))^{-1}}_\kappa\le2.
 \label{na-chart-weighted}
\end{equation}
These are respectively the $B$ derivative of the chart and its
$\xi$ derivative restricted to $\ker Q$. The global bounded-field
norms have the corresponding bounds with $k$. Holomorphy of the
weighted inverse and Cauchy on smaller parameter balls give
factorial higher derivatives as well.
\end{proof}

If $\xi$ additionally satisfies the fixed flat gauge equation
$RD^*\xi=0$, the identity $RD^*H=0$ and
\eqref{na-fixed-point} give $RD^*A=0$. On finite regulators this
is immediate by continuity. On the infinite lattice it holds on the
domain of the fixed gauge operator, or its bounded-cell extension;
no new uniform bounded-field estimate for the order-zero gauge
projection is asserted.

\subsection{Scope of the result}
\label{na-scope}

This proves a refinement-uniform complex ball for the complete
corrected nonlinear iterate, its physical Hilbert differential
bounds and an exact analytic straightening of its constraint along
the flat minimizing-extension range. The result is internally
audited, not externally certified.

The coordinate is the fixed flat-reference one in
\eqref{na-coordinate}. At nonzero fields its ordinary derivative
has not here been equated with a source operator defined through
left multiplication, new background-dependent endpoint corrections
or a moving gauge. That identification requires the actual
coordinate Jacobians and reference-background transformation.
The selected chart field is not by itself an action minimizer.
Notebook 14 supplies the reference-coordinate correction and a
nonflat energy inverse. Notebook 15 subsequently solves the
nonlinear stationarity equation in this chart, proves a strict
fixed-slice minimum, and transfers it by actual gauges to the
ordinary homogeneous constraint. Those later results remain
restricted to a small global physical $C^{1,\alpha}$ ball.
Retained or mixed-scale geometry and an interacting continuum
limit remain separate tasks.

\begin{thebibliography}{9}
\bibitem{repo}
\href{https://github.com/ethanwsussman/Yang--Mills-LLM-experiment/tree/e392ef45013a581664bab76b90867170b30fe10e}
{Pinned Yang--Mills-LLM-experiment account}, commit
\texttt{e392ef45013a581664bab76b90867170b30fe10e}.
The corrected nonlinear composition is Raw I:24320--24325;
its coordinate and flat differential are I:24344--24369;
the endpoint-correction bridge is I:47076--47087.
\bibitem{previous}
Companion research proofs: Notebook 08, the source-compatible flat
minimizing extension; Notebook 11, the full flat constraint-current
family; Notebook 12, pointwise regularity on the flat current domain.
The original CMP98 full paper has not been independently reverified.
\end{thebibliography}
\end{document}
